% Part: proof-theory
% Chapter: proof-search
% Section: completeness

\documentclass[../../../include/open-logic-section]{subfiles}

\begin{document}

\olfileid{pt}{ps}{cpl}

\olsection{संपूर्णता}

आता सिद्धता-शोध अल्गोरिदम विश्वासार्ह आहे हे दाखवू:
तो थांबला तरी किंवा कधीच समाप्त झाला नाही तरी,
सुरुवातीच्या क्रमवर्ती~$\Gamma \Sequent \Delta$ ची
!!{proof} नसते. यासाठी अशी !!a{structure}~$\Struct M$
बांधू की प्रत्येक~$!A \in \Gamma$ साठी~$\Sat{M}{!A}$
आणि प्रत्येक~$!A \in \Delta$ साठी~$\Sat/{M}{!A}$.
म्हणजे~$\Gamma$ मधील सर्व !!{formula}s सत्य,
तर~$\Delta$ मधील सर्व !!{formula}s हे~$\Struct M$
मध्ये असत्य आहेत. त्यामुळे~$\Gamma \Sequent \Delta$
वैध नाही. \Log{G3c} निर्दोष असल्याने
$\Gamma \Sequent \Delta$ ची !!{proof} नाही.

!!{formula}s ची जोडी~$\Theta, \Xi$ आणि
!!{constant}s चा संच~$C$ यांपासून~$\Struct{M}$ बांधतो.
हे~$\Theta, \Xi$ आणि~$C$ आपल्याला अयशस्वी
(थांबलेल्या किंवा समाप्त न होणाऱ्या) सिद्धता-शोधाच्या
\emph{अपयश-शाखे}तून मिळतात. अपयश-शाखा म्हणजे
क्रमवर्तींचा~$\Pi_n \Sequent \Lambda_n$ अनुक्रम
आणि~!!{constant}s चे संच~$C_n$; येथे
$\Pi_0 \Sequent \Lambda_0$ ही सुरुवातीची
अंतिम क्रमवर्ती~$\Gamma \Sequent \Delta$ असते,
आणि~$\Pi_{n+1} \Sequent \Lambda_{n+1}$
हे~$n+1$ व्या टप्प्यावर~$\Pi_n \Sequent \Lambda_n$
पासून निर्माण झालेले असे आधारविधान असते की
अल्गोरिदमला त्याची !!a{proof} सापडत नाही.
प्रत्येक~$n$ साठी असे आधारविधान असलेच पाहिजे;
नसते तर अल्गोरिदमला !!a{proof} सापडली असती.
~$C_n$ हा~$\Pi_n \Sequent \Lambda_n$
शी संबंधित !!{constant}s चा संच घेऊ.

अल्गोरिदम फक्त अणुरूप !!{formula}s असलेल्या
सर्वांत वरच्या क्रमवर्तीवर थांबला, तर त्या
क्रमवर्तीवर संपणारी शाखा ही अपयश-शाखा आहे.
अल्गोरिदम कधीच समाप्त झाला नाही, तर तो वाढत्या
मोठ्या वृक्षांची निर्मिती करत राहतो. हे वृक्ष
वाढून एक अनंत वृक्ष तयार होतो असे समजू शकतो
(अल्गोरिदम अनंत टप्पे चालल्यानंतर जो वृक्ष असेल तो).
हा अनंत वृक्ष सांत-शाखित असतो: प्रत्येक शिखराला
फक्त एक किंवा दोन बालशिखरे असतात---म्हणजे त्याच्या
तत्काळ वरच्या क्रमवर्ती. क्योनिग (K\H{o}nig)
यांच्या उपप्रमेयानुसार,
प्रत्येक अनंत सांत-शाखित वृक्षाला अनंत शाखा असते.
सिद्धता-शोध कायम चालू राहिल्यास अशी अनंत शाखा
अपयश-शाखा ठरते.

~$\Pi_n \Sequent \Lambda_n$ आणि~$C_n$
यांचा अनुक्रम अपयश-शाखा असेल, तर
$\Theta = \bigcup_n \Pi_n$ आणि~$\Xi = \bigcup_n \Lambda_n$
घेऊ (!!{formula}s वरील निर्देशांक आणि त्यांच्या
अनेक आस्थिती काढून); तसेच~$C = \bigcup_n C_n$ घेऊ.

आता~$\Struct{M}$ परिभाषित करू.
\begin{enumerate}
\item $\Domain{M}$ हा~$C$ मधील स्थिरांक आणि
$\Gamma \Sequent \Delta$ मधील !!{function}s
(असतील तर) वापरून, पण !!{variable}s न वापरता,
बनवता येणाऱ्या पदांचा संच आहे.
\item $c \in C$ असेल, तर~$\Assign{c}{M} = c$.
\item $f$ हे~$\Gamma \Sequent \Delta$ मधील
$m$-स्थानी !!{function} असेल आणि~$t_1$,
\dots, $t_m \in \Domain{M}$ असतील, तर
$\Assign{f}{M}(t_1, \dots, t_m) = f(t_1, \dots, t_m)$.
\item $R$ हे~$\Gamma \Sequent \Delta$ मधील
$m$-स्थानी !!{predicate} असेल, तर
$\Assign{R}{M} = \Setabs{\tuple{t_1, \dots, t_m} \in
\Domain{M}^m}{R(t_1, \dots, t_m) \in \Theta}$.
\end{enumerate}
~$\Struct{M}$ हे \emph{पद-प्रतिमान} आहे, कारण
त्याचे विश्व पदांचा संच आहे आणि !!{constant}s व
!!{function}s यांचा अर्थ असा लावला आहे की
$\Value{t}{M} = t$ मिळते. ~$\Theta$ मधील अणुरूप
!!{formula}s वापरून~$\Assign{R}{M}$ अशी परिभाषित
केली आहे की~$\Sat{M}{R(t_1, \dots, t_m)}$ तेव्हाच
आणि तेव्हाच जेव्हा~$R(t_1, \dots, t_m) \in \Theta$.

\begin{lem}\ollabel{lem:truth}
प्रत्येक~$!A \in \Theta \cup \Xi$ साठी:
\begin{enumerate}
  \item $!A \in \Theta$ असेल, तर~$\Sat{M}{!A}$.
  \item $!A \in \Xi$ असेल, तर~$\Sat/{M}{!A}$.
\end{enumerate}
\end{lem}

\begin{proof}
~$\depth{!A}$ वर विगमन करू.

प्रथम~$!A$ अणुरूप आहे असे मानू;
म्हणजे~$!A \ident \lfalse$ किंवा
$!A \ident R(t_1, \dots, t_m)$.

~$\Pi_n \Sequent \Lambda_n$ ही अपयश-शाखा
असल्याने कोणत्याही~$\Pi_n$ मध्ये~$\lfalse$
असू शकत नाही (नाहीतर~$\Pi_n \Sequent \Lambda_n$
स्वयंसिद्धक ठरेल). ~$\Struct{M}$ च्या व्याख्येनुसार
$\Sat{M}{R(t_1, \dots, t_m)}$ तेव्हाच आणि तेव्हाच
जेव्हा~$R(t_1, \dots, t_m) \in \Theta$.
दोन्ही मिळून: $!A \in \Theta$ असेल, तर
$\Sat{M}{!A}$.

~$!A$ अणुरूप असून~$!A \in \Pi_n$ असेल,
तर~$!A \in \Pi_{n+1}$; आणि~$!A \in \Lambda_n$
असेल, तर~$!A \in \Lambda_{n+1}$.
(सिद्धता-शोध अल्गोरिदम पुढील क्रमवर्ती कशा
बनवतो यावरून हे मिळते.) म्हणून~$!A \in \Theta \cap \Xi$
असेल, तर एखाद्या~$n$ साठी~$!A \in \Pi_n \cap
\Lambda_n$. मग~$\Pi_n \Sequent \Lambda_n$
स्वयंसिद्धक ठरते; पण अपयश-शाखेत स्वयंसिद्धके
असू शकत नाहीत. म्हणून अणुरूप~$!A$ साठी
$!A \notin \Theta \cap \Xi$; परिणामी~$!A \in \Xi$
असेल तर~$!A \notin \Theta$. त्यामुळे~$\Struct{M}$
च्या व्याख्येनुसार~$!A \in \Xi$ असताना
$\Sat/{M}{!A}$.

विगमन-टप्प्यासाठी,~$!A$ पेक्षा कमी
!!{depth} असलेल्या सर्व !!{formula}s साठी
(1) आणि~(2) खरे आहेत असे मानू. ~$!A$ साठी
(1) व~(2) सिद्ध करण्यासाठी प्रसंग वेगळे पाहू.

प्रथम,~$\Pi_n \Sequent \Lambda_n$ मध्ये
$!A^i$ आल्यास ते~$\Pi_{n+1} \Sequent \Lambda_{n+1}$
मध्येही~$!A^i$ म्हणून येते, जोपर्यंत~$i$ हा~$\Pi_n \Sequent \Lambda_n$
मधील सर्वांत लहान निर्देशांक नसतो. प्रत्येक
टप्प्यावर जोडलेल्या सर्वांत वरच्या क्रमवर्तींत
सर्वांत लहान निर्देशांक वाढतो. म्हणून अखेरीस
$!A^i$ असलेली अशी~$\Pi_n \Sequent \Lambda_n$
मिळते की~$i$ हा सर्वांत लहान निर्देशांक असतो.
~$n+1$ व्या टप्प्यावर अल्गोरिदम~$!A^i$ चे
न्यूनीकरण करतो. दुसऱ्या शब्दांत,
$!A \in \Theta$ असेल, तर एखाद्या~$n$ साठी
$\Pi_n = \Pi_n', !A^i$ आणि~$i$ सर्वांत
लहान निर्देशांक असतो. तसेच~$!A \in \Xi$
असेल, तर एखाद्या~$n$ साठी
$\Lambda_n = \Lambda_n', !A^i$ आणि~$i$
सर्वांत लहान निर्देशांक असतो.

\begin{enumerate}
  \item $!A \in \Theta$ आणि~$!A \ident !B \land !C$
  असेल, तर एखाद्या~$n$ साठी
  $\Pi_n = \Pi_n', (!B \land !C)^i$.
  ~$\Pi_{n+1} \Sequent \Lambda_{n+1}$ हे
  \LeftR{\land} चे अनुरूप आधारविधान आहे;
  म्हणजे~$\Pi_{n+1} = \Pi_n', !B^k, !C^{k+1}$.
  त्यामुळे~$!B, !C \in \Theta$. विगमन गृहीतकानुसार
  $\Sat{M}{!B}$ आणि~$\Sat{M}{!C}$.
  ~$\Sat{M}{}$ च्या व्याख्येवरून
  $\Sat{M}{!B \land !C}$.

  \item $!A \in \Xi$ आणि~$!A \ident !B \land !C$
  असेल, तर एखाद्या~$n$ साठी
  $\Lambda_n = \Lambda_n', !B \land !C$.
  ~$\Pi_{n+1} \Sequent \Lambda_{n+1}$ हे
  निष्कर्ष~$\Pi_n \Sequent \Lambda'_n, !B \land !C$
  असलेल्या~\RightR{\land} च्या दोन आधारविधानांपैकी
  एक आहे; म्हणजे~$\Lambda_{n+1} = \Lambda_n', !B^k$
  किंवा~$\Lambda_{n+1} = \Lambda_n', !C^k$.
  त्यामुळे~$!B \in \Xi$ किंवा~$!C \in \Xi$.
  विगमन गृहीतकानुसार~$\Sat/{M}{!B}$ किंवा
  $\Sat/{M}{!C}$. ~$\Sat{M}{}$ च्या व्याख्येवरून
  $\Sat/{M}{!B \land !C}$.

  \item $!A \ident \lnot !B$,
  $!A \ident !B \lor !C$ आणि~$!A \ident !B \lif !C$
  हे प्रसंगही याच प्रकारे हाताळता येतात;
  ते सरावासाठी ठेवले आहेत.

  \item $!A \in \Theta$ आणि
  $!A \ident \lexists[x][!B(x)]$ असेल,
  तर एखाद्या~$n$ साठी
  $\Pi_n = \Pi_n', \lexists[x][!B(x)]^i$.
  ~$\Pi_{n+1} \Sequent \Lambda_{n+1}$ हे
  \LeftR{\lexists} चे अनुरूप आधारविधान आहे;
  म्हणजे एखाद्या~$c$ साठी
  $\Pi_{n+1} = \Pi_n', !B(c)^k$.
  त्यामुळे~$!B(c) \in \Theta$ आणि~$c \in C$.
  विगमन गृहीतकानुसार~$\Sat{M}{!B(c)}$.
  ~$\Sat{M}{}$ च्या व्याख्येवरून
  $\Sat{M}{\lexists[x][!B(x)]}$.

  \item $!A \in \Xi$ आणि
  $!A \ident \lexists[x][!B(x)]$ असेल,
  तर एखाद्या~$n$ साठी
  $\Lambda_n = \Lambda_n', \lexists[x][!B(x)]^i$.
  प्रत्येक वेळी~$\lexists[x][!B(x)]$ चे
  न्यूनीकरण करताना~$\lexists[x][!B(x)]$
  आधारविधानाच्या उजव्या
  बाजूस नव्या निर्देशांकासह राहते. म्हणून
  अपयश-शाखेत ते उजव्या बाजूस असेल, तर
  $\lexists[x][!B(x)]$ चे अनंत वेळा
  न्यूनीकरण होते. प्रत्येक~$t \in \Domain{M}$
  हे अखेरीस अपयश-शाखेत~``$C_n$ मधीलच स्थिरांक
  असलेले आणि उजवीकडील~$\lexists[x][!B(x)]$
  च्या आधीच्या कोणत्याही न्यूनीकरणात न वापरलेले
  पहिले पद'' ठरते. त्यामुळे प्रत्येक
  $t \in \Domain{M}$ साठी एखाद्या~$n$ ला
  $!B(t) \in \Lambda_n$, आणि म्हणून~$!B(t) \in \Xi$.
  विगमन गृहीतकानुसार प्रत्येक~$t \in \Domain{M}$
  साठी~$\Sat/{M}{!B(t)}$. ~$\Sat{M}{}$
  च्या व्याख्येवरून~$\Sat/{M}{\lexists[x][!B(x)]}$.

  \item $!A \ident \lforall[x][!B(x)]$
  हे प्रसंगही याच प्रकारे हाताळता येतात;
  ते सरावासाठी ठेवले आहेत.
\end{enumerate}
\end{proof}

\begin{cor}\ollabel{cor:complete}
~\Log{G3c} साठीचा सिद्धता-शोध अल्गोरिदम संपूर्ण
आहे: तो थांबला किंवा कधीच समाप्त झाला नाही,
तर सुरुवातीची क्रमवर्ती~$\Gamma \Sequent \Delta$
वैध नसते आणि म्हणून तिची !!{proof} नसते.
\end{cor}

\begin{proof}
  अल्गोरिदम थांबला किंवा कधीच समाप्त झाला नाही,
  तर अशी अपयश-शाखा असते की तिच्यापासून
  $\Struct{M}$ परिभाषित करता येते.
  \olref{lem:truth} नुसार प्रत्येक~$!A \in \Gamma$
  साठी~$\Sat{M}{!A}$ आणि प्रत्येक~$!A \in \Delta$
  साठी~$\Sat/{M}{!A}$. (लक्षात घ्या की
  $\Pi_0 = \Gamma$ आणि~$\Lambda_0 = \Delta$;
  म्हणून~$\Gamma \subseteq \Theta$ आणि
  $\Delta \subseteq \Xi$.) त्यामुळे
  $\Gamma \Sequent \Delta$ वैध नाही.
  ~\Log{G3c} च्या निर्दोषतेवरून
  $\Gamma \Sequent \Delta$ ची !!{proof} नाही.
\end{proof}

\begin{cor}
~\Log{G3c} संपूर्ण आहे:
$\Gamma \Sequent \Delta$ वैध असेल,
तर~$\Log{G3c} \Proves \Gamma \Sequent \Delta$.
\end{cor}

\begin{proof}
  व्यत्यस्त विधान सिद्ध करू.
  $\Log{G3c} \Proves/ \Gamma \Sequent \Delta$
  असे मानू. मग शोधण्याजोगी !!{proof} नसल्याने
  सिद्धता-शोध अल्गोरिदम थांबला पाहिजे किंवा
  कधीच समाप्त झाला नाही पाहिजे.
  \olref{cor:complete} नुसार
  $\Gamma \Sequent \Delta$ वैध नाही.
\end{proof}

\end{document}
